\documentclass[10pt, letterpaper]{article}

% Packages:
\usepackage[
    ignoreheadfoot,
    top=2 cm,
    bottom=2 cm,
    left=2 cm,
    right=2 cm,
    footskip=1.0 cm,
]{geometry}
\usepackage{titlesec}
\usepackage{tabularx}
\usepackage{array}
\usepackage[dvipsnames]{xcolor}
\definecolor{primaryColor}{RGB}{0, 79, 144}
\usepackage{enumitem}
\usepackage{fontawesome5}
\usepackage{amsmath}
\usepackage[
    pdftitle={Rudra Mani Upadhyay - Resume},
    pdfauthor={Rudra Mani Upadhyay},
    pdfcreator={Rudra},
    colorlinks=true,
    urlcolor=primaryColor
]{hyperref}
\usepackage[pscoord]{eso-pic}
\usepackage{calc}
\usepackage{bookmark}
\usepackage{lastpage}
\usepackage{changepage}
\usepackage{paracol}
\usepackage{ifthen}
\usepackage{needspace}
\usepackage{iftex}

% PDF Readability:
\ifPDFTeX
    \input{glyphtounicode}
    \pdfgentounicode=1
    \usepackage[utf8]{inputenc}
    \usepackage{lmodern}
\fi

% Page Setup:
\AtBeginEnvironment{adjustwidth}{\partopsep0pt}
\pagestyle{empty}
\setcounter{secnumdepth}{0}
\setlength{\parindent}{0pt}
\setlength{\topskip}{0pt}
\setlength{\columnsep}{0cm}
\makeatletter
\let\ps@customFooterStyle\ps@plain
\patchcmd{\ps@customFooterStyle}{\thepage}{
    \color{gray}\textit{\small Rudra - Page \thepage{} of \pageref*{LastPage}}
}{}{}
\makeatother
\pagestyle{customFooterStyle}

\titleformat{\section}{\needspace{4\baselineskip}\bfseries\large}{}{0pt}{}[\vspace{1pt}\titlerule]
\titlespacing{\section}{-1pt}{0.3 cm}{0.2 cm}
\renewcommand\labelitemi{$\circ$}

% Custom environments:
\newenvironment{highlights}{
    \begin{itemize}[
        topsep=0.10 cm,
        parsep=0.10 cm,
        partopsep=0pt,
        itemsep=0pt,
        leftmargin=0.4 cm + 10pt
    ]
}{
    \end{itemize}
}

\newenvironment{highlightsforbulletentries}{
    \begin{itemize}[
        topsep=0.10 cm,
        parsep=0.10 cm,
        partopsep=0pt,
        itemsep=0pt,
        leftmargin=10pt
    ]
}{
    \end{itemize}
}

\newenvironment{onecolentry}{
    \begin{adjustwidth}{0.2 cm + 0.00001 cm}{0.2 cm + 0.00001 cm}
}{
    \end{adjustwidth}
}

\newenvironment{twocolentry}[2][]{
    \onecolentry
    \def\secondColumn{#2}
    \setcolumnwidth{\fill, 4.5 cm}
    \begin{paracol}{2}
}{
    \switchcolumn \raggedleft \secondColumn
    \end{paracol}
    \endonecolentry
}

\newenvironment{header}{
    \setlength{\topsep}{0pt}\par\kern\topsep\centering\linespread{1.5}
}{
    \par\kern\topsep
}

\newcommand{\placelastupdatedtext}{
  \AddToShipoutPictureFG*{
    \put(
        \LenToUnit{\paperwidth-2 cm-0.2 cm+0.05cm},
        \LenToUnit{\paperheight-1.0 cm}
    ){\vtop{{\null}\makebox[0pt][c]{
        \small\color{gray}\textit{Last updated in October 2025}\hspace{\widthof{Last updated in October 2025}}
    }}}%
  }%
}

\let\hrefWithoutArrow\href
\renewcommand{\href}[2]{\hrefWithoutArrow{#1}{\ifthenelse{\equal{#2}{}}{ }{#2 }\raisebox{.15ex}{\footnotesize \faExternalLink*}}}

\begin{document}
\newcommand{\AND}{\unskip\cleaders\copy\ANDbox\hskip\wd\ANDbox\ignorespaces}
\newsavebox\ANDbox
\sbox\ANDbox{}

\placelastupdatedtext
\begin{header}
    \textbf{\fontsize{24 pt}{24 pt}\selectfont Rudra Mani Upadhyay}

    \vspace{0.3 cm}

    \normalsize
    \mbox{{\color{black}\footnotesize\faMapMarker*}\hspace*{0.13cm}New Delhi, India}%
    \kern 0.25 cm%
    \AND%
    \kern 0.25 cm%
    \mbox{\hrefWithoutArrow{mailto:rudramaniupadhyay651@gmail.com}{\color{black}{\footnotesize\faEnvelope[regular]}\hspace*{0.13cm}rudramaniupadhyay651@gmail.com}}%
    \kern 0.25 cm%
    \AND%
    \kern 0.25 cm%
    \mbox{\hrefWithoutArrow{tel:+919465784195}{\color{black}{\footnotesize\faPhone*}\hspace*{0.13cm}+91 9465784195}}%
    \kern 0.25 cm%
    \AND%
    \kern 0.25 cm%
    \mbox{\hrefWithoutArrow{https://www.linkedin.com/in/rudra-mani-upadhyay-47832b289}{\color{black}{\footnotesize\faLinkedinIn}\hspace*{0.13cm}LinkedIn}}%
    \kern 0.25 cm%
    \AND%
    \kern 0.25 cm%
    \mbox{\hrefWithoutArrow{https://github.com/Rudramani1}{\color{black}{\footnotesize\faGithub}\hspace*{0.13cm}GitHub}}%
\end{header}

% --- Ordered Sections ---
\section{Professional Summary}
\begin{onecolentry}
    Detail-oriented and analytical B.Tech student specializing in Artificial Intelligence and Machine Learning with hands-on experience in Python, Flask, and full-stack web development. Strong foundation in algorithms, data analysis, and database systems. Experienced in building and deploying AI-driven applications using modern frameworks and APIs. Seeking opportunities to apply technical and problem-solving skills to real-world software and AI challenges.
\end{onecolentry}

\section{Technical Skills}
\begin{onecolentry}
    \textbf{Programming Languages:} Python, Java, C++, JavaScript, SQL, R, Bash

    \textbf{Frameworks \& Libraries:} Flask, Django, TensorFlow, NumPy, Pandas, Scikit-learn, Bootstrap, React.js, Node.js, Next.js

    \textbf{Databases:} MySQL, Firebase, SQLite

    \textbf{Tools \& Platforms:} Git, Docker, VS Code, Jupyter Notebook, Linux, Google Cloud, Postman

    \textbf{Soft Skills:} Team Collaboration, Adaptability, Time Management, Communication, Problem Solving, Logical Reasoning
\end{onecolentry}

\section{Experience}

\begin{twocolentry}{
    \textit{May 2025 -- July 2025}
}
    \textbf{FOSSEE Summer Fellowship} \\
    \textit{IIT Bombay | Remote}
\end{twocolentry}

\begin{onecolentry}
    \begin{highlights}
        \item Developed an AI-powered \textbf{eSim ChatBot} using Python and Flask for circuit simulation support.
        \item Integrated \textbf{Large Language Models (LLMs)} with domain-specific datasets to improve contextual response accuracy by 25\%.
        \item Deployed the chatbot in a scalable environment and contributed to FOSSEE’s open-source ecosystem.
    \end{highlights}
\end{onecolentry}

\begin{twocolentry}{
    \textit{June 2025 -- July 2025}
}
    \textbf{Astronomy and Astrophysics Summer School} \\
    \textit{India Space Academy}
\end{twocolentry}

\begin{onecolentry}
    \begin{highlights}
        \item Completed \textbf{120 hours of intensive training and project work} in Data-Driven Astronomy.
        \item Applied computational methods to analyze astronomical datasets and astrophysical phenomena using Python.
        \item Strengthened skills in scientific inquiry, data visualization, and research methodology.
    \end{highlights}
\end{onecolentry}

\section{Projects}

\begin{twocolentry}{
    \textit{\href{https://github.com/Rudramani1/Youtube-Summarizer}{GitHub Repository}}
}
    \textbf{YouTube Video Summarizer}
\end{twocolentry}
\begin{onecolentry}
    \begin{highlights}
        \item Built a Flask web app that extracts and summarizes YouTube video transcripts using \textbf{Together AI API} and NLP techniques.
        \item Implemented API integration, web scraping, and text summarization for efficient content generation.
        \item Tools: Flask, HTML, CSS, JavaScript, YouTube Transcript API, NLP.
    \end{highlights}
\end{onecolentry}

\vspace{0.2 cm}

\begin{twocolentry}{
    \textit{\href{https://github.com/Rudramani1/Pneumonia-Analysis}{GitHub Repository}}
}
    \textbf{Pneumonia Detection \& Awareness Platform}
\end{twocolentry}
\begin{onecolentry}
    \begin{highlights}
        \item Developed a web app for pneumonia awareness and X-ray analysis using \textbf{Teachable Machine (TensorFlow.js)}.
        \item Implemented AI-based image classification and interactive awareness dashboard.
        \item Tools: TensorFlow.js, Bootstrap, HTML, CSS, JavaScript.
    \end{highlights}
\end{onecolentry}

\section{Certifications \& Achievements}
\begin{onecolentry}
    \begin{highlights}
        \item \textbf{NPTEL:} Java Programming (IIT) \quad | \quad \textbf{Infosys Certification:} Data Structures \& Algorithms in Python
        \item Certified in Python Programming, SQL, HTML \& CSS, and Machine Learning Essentials.
        \item \textbf{NCAT (National Creativity Aptitude Test 2025):} Rank 12
        \item \textbf{INNOQUEST Hackathon:} 3rd Place \quad | \quad \textbf{Smart AI Hackathon:} 3rd Place
    \end{highlights}
\end{onecolentry}

\section{Education}

\begin{twocolentry}{
    \textit{Sept 2023 -- May 2027}
}
    \textbf{Dronacharya College of Engineering, Gurugram, Haryana} \\
    \textit{B.Tech in Computer Science (Artificial Intelligence \& Machine Learning)}
\end{twocolentry}

\begin{onecolentry}
    \begin{highlights}
        \item Relevant Coursework: Data Structures, Algorithms, Operating Systems, Database Management Systems (DBMS), Computer Networks, Data Science, Machine Learning.
    \end{highlights}
\end{onecolentry}

\begin{twocolentry}{
    \textit{2021 -- 2022}
}
    \textbf{Kendriya Vidyalaya Rangpuri AAI, New Delhi} \\
    \textit{Senior Secondary (PCM + CS) – 79.8\%}
\end{twocolentry}

\begin{twocolentry}{
    \textit{2019 -- 2020}
}
    \textbf{Kendriya Vidyalaya Rangpuri AAI, New Delhi} \\
    \textit{Secondary (10\textsuperscript{th}) – 82\%}
\end{twocolentry}

\end{document}